\subproblem{Переменный (?) ток}{3.5}

\pic7R{0pt}{0.3}{figures/1b.1}

В цепи, показанной на рисунке справа, последовательно соединённые конденсатор ёмкостью $C=\qty{4.7}{\mF}$ и катушка индуктивностью $L=\qty{5.4}{\H}$ подключены к ключу $K$, который в обоих пунктах подключается к источнику постоянного напряжения $U_0=\qty{10}{\V}$ на очень короткое время $\tau=\qty5{\ms}$, затем в течение времени $T=\qty{250}{\ms}$ (считайте, что $T\gg\tau$) переключен на короткое замыкание, т.е. крайне левое вертикальное положение, а потом снова подключается к источнику, и так далее. В момент $t=0$ тока в цепи не было в обоих пунктах.

\begin{task}[type=subproblem]
    \item Пусть конденсатор был разряжен, а ключ находился в крайнем левом вертикальном положении. В момент $t=0$ ключ замкнулся на $U_0$ и продолжил переключаться так, как описано выше. Каким станет энергия системы в момент $t=3T+\tau$ (т.е. $t=3T$ сразу после переключения $K$ на крайнее левое положение)? Чему был максимальный заряд конденсатора $q_0$ за промежуток $0\leq t \leq 3T+\tau$?
    \item Допустим, в момент $t=0$ на конденсаторе был неизвестный заряд $q_*>0$ (полярность указана на рисунке), а ключ неподвижно находился в крайнем левом вертикальном положении. Затем, в некоторый момент $t=t_*>0$, ключ замкнули на источник $U_0$ и продолжили переключать как описано выше. Оказалось, полная энергия системы осталась неизменной при переключении ключа $K$. При каких $q_*$ и $t_*$ такое возможно?
\end{task}
